\section{Justificación}

El fortalecimiento del sector cafetero colombiano en el segmento de cafés de 
especialidad requiere tecnologías que mejoren la calidad del producto, la eficiencia 
de los procesos y la sostenibilidad del sistema productivo. En este contexto, el secado 
constituye una etapa crítica de la cadena de valor por su influencia directa sobre las 
propiedades físicas, químicas y sensoriales del grano y su valor comercial 
\parencite{abreu_influence_2025}. Esto cobra especial relevancia ante el crecimiento 
del segmento de cafés de especialidad, que avanza a tasas cercanas al 15\% anual frente 
al 2\% del café tipo \textit{commodity}, posicionando al secado como una operación 
estratégica para la preservación del valor agregado \parencite{barbosa_sensory_2019}.

Desde una perspectiva social, la caficultura colombiana involucra aproximadamente 
540000 familias productoras, de las cuales la gran mayoría son pequeños caficultores 
con un promedio de dos hectáreas sembradas por finca, generando cerca de 2.5 millones 
de empleos directos e indirectos en el país \parencite{betancourt_agricultural_2024, tabarquino_munoz_cadena_2025}. La mayor 
concentración de productores se ubica en la región Andina, donde departamentos como 
Antioquia (16.6\%), Cauca (11.4\%), Huila (13.7\%), Tolima (10.7\%) y Nariño representan 
conjuntamente el 54.8\% del área cultivada nacional \parencite{betancourt_agricultural_2024, tabarquino_munoz_cadena_2025, navarro_munoz_sector_2024}, 
siendo el departamento del Valle del Cauca un escenario estratégico para la 
validación y escalamiento de tecnologías de beneficio centralizado. Estudios previos evidencian una correlación directa 
entre mayores rendimientos productivos y menores índices de pobreza multidimensional 
en zonas rurales cafeteras \parencite{betancourt_agricultural_2024}, lo que sugiere que la 
tecnificación del secado no solo impacta la calidad del producto, sino también la 
estabilidad del ingreso familiar y las condiciones de vida de las comunidades 
productoras.

En materia de equidad de género, las mujeres representan el 26\% de los 
caficultores \parencite{navarro_munoz_sector_2024}, y la reducción de la carga física 
asociada al secado solar mediante sistemas mecanizados contribuye a su inclusión 
efectiva en la cadena productiva \parencite{tabarquino_munoz_cadena_2025, betancourt_agricultural_2024}.

En términos económicos, el mercado internacional de café atraviesa una coyuntura de 
contracción en volumen con sostenimiento de precios, lo que amplifica la importancia 
estratégica de la calidad diferenciada. Según el informe técnico de la Federación 
Nacional de Cafeteros de febrero de 2026, Colombia exportó 846588 sacos de 60\,kg, 
registrando una caída del 28.5\% frente al mismo período de 2025, con un valor FOB 
estimado de USD 326.8 millones; no obstante, el precio medio se incrementó un 11.9\%, 
situándose en 291.8\,\textcent/lb \parencite{direccion_de_investigaciones_economicas_fnc_informe_2026}. En este 
escenario, el segmento de cafés diferenciados —cuya demanda crece al 15\% anual frente 
al 2\% del café \textit{commodity} \parencite{barbosa_sensory_2019} y cuyos precios 
pueden superar hasta diez veces el valor del café convencional 
\parencite{tabarquino_munoz_cadena_2025}— emerge como el vector de mayor potencial para 
ampliar el valor exportado. Sin embargo, procesos de secado deficientes pueden devaluar 
el precio del grano entre un 10\% y un 20\% \parencite{ferreira_bernardes_pereira_influence_2020}, 
mientras que el uso de aire deshumidificado controlado puede incrementar su valor 
comercial en más del 12\% \parencite{konopatzki_price_2019}.

A esto se suma que el 
secado representa más del 60\% del consumo energético total del procesamiento 
\parencite{coradi_development_2021, cruz-ospina_improving_2026}, y que las cosechas 
colombianas coinciden con épocas de alta pluviosidad y humedad relativa superior al 
70\%, condición que agrava los riesgos del secado solar tradicional 
\parencite{largoavila_changes_2023, osorio_hernandez_bioclimatic_2018}. Ante estas limitaciones, la implementación de 
plantas de beneficio centralizadas bajo el esquema CERPER ofrece un modelo de acceso 
colectivo que democratiza el uso de tecnologías avanzadas de secado entre productores 
de menor escala \parencite{khathir_drying_2024}.

Técnicamente, los sistemas convencionales basados en control de temperatura pueden 
inducir daños térmicos superiores a los 40--45\,$^\circ$C, afectando la integridad 
celular y la calidad sensorial del grano \parencite{borem_quality_2018, 
oliveira_quality_2018, alves_physiological_2017}. En contraste, el control de la 
humedad relativa del aire incrementa la tasa de secado sin comprometer la calidad, al 
aumentar la difusividad efectiva de la humedad \parencite{phitakwinai_thinlayer_2019, 
andrade_mathematical_2019, borem_quality_2018}. El secado asistido por bomba de calor 
ha demostrado mejorar significativamente la eficiencia energética del proceso 
\parencite{fauzi_influence_2024}, aunque persiste una brecha en su implementación y 
validación a escala industrial. Esta brecha es especialmente crítica en cafés 
diferenciados como \textit{natural} y \textit{honey}, donde una mayor carga de humedad 
y sensibilidad al proceso puede derivar en defectos sensoriales, crecimiento microbiano 
y pérdida de valor comercial \parencite{borem_quality_2018, parra-coronado_preliminary_2020}.

Desde la perspectiva ambiental, la bomba de calor permite recuperar el calor latente 
del aire de salida, reduciendo el consumo energético y la dependencia de combustibles 
fósiles \parencite{fauzi_influence_2024, guevara-sanchez_effect_2019, 
cruz-ospina_improving_2026}. La integración de sistemas fotovoltaicos posibilita 
configuraciones híbridas que reducen la huella de carbono del proceso 
\parencite{prada_efectividad_2019}, en línea con los objetivos de la Misión de Energía 
Sostenible de la Política de Investigación e Innovación Orientada por Misiones (PIIOM) 
del Ministerio de Ciencias.

En el plano institucional, el Estado colombiano ha reconocido la importancia estratégica 
del sector mediante el CONPES 4052 de 2021, denominado ``Política para la Sostenibilidad 
de la Caficultura Colombiana'', orientado a estabilizar los ingresos cafeteros e impulsar 
la comercialización diferenciada, en articulación con 
el Programa de Cafés Especiales de la Federación Nacional de Cafeteros (FNC), activo 
desde 1986, y con centros de investigación como CENICAFÉ \parencite{tabarquino_munoz_cadena_2025}. 
Asimismo, el control preciso de humedad y temperatura durante el secado resulta 
determinante para cumplir con los estándares de calidad exigidos por las Denominaciones 
de Origen y el sello Café de Colombia, instrumentos de competitividad que diferencian 
el producto colombiano en mercados internacionales \parencite{tabarquino_munoz_cadena_2025, abreu_influence_2025, jakkaew_data-driven_2024}.

Finalmente, la investigación se articula con estrategias del sector cafetero colombiano 
como las plantas de beneficio centralizadas y el modelo de compra de cereza en su equivalente en pergamino seco (CERPER), donde tecnologías de 
secado controladas y escalables son fundamentales para garantizar una adecuada operación. Se identifica así una brecha 
tecnológica en el desarrollo de sistemas que integren control preciso de la humedad 
relativa, eficiencia energética y escalabilidad \parencite{khathir_drying_2024, 
konopatzki_coffee_2022, malik_microstructure_2020}, justificando el escalado de una 
tecnología de secado asistido por bomba de calor para estandarizar procesos de cafés 
diferenciados y mejorar la competitividad del sector cafetero colombiano.

% (pegar en el cuerpo del documento, donde quieras que aparezca)

\begin{figure}[htbp]
\centering
\resizebox{\textwidth}{!}{%
\begin{tikzpicture}[
    font=\sffamily,
    every node/.style={align=center, text=textDark},
]

% --- Título ---
\node[font=\bfseries\LARGE, text=red1] at (0, 2.0)
    {Diagrama de Reloj de Arena};
\node[font=\large\bfseries, text=purple1] at (0, 1.0)
    {Justificación de la Investigación};

% ==================================================================
%  PARTE SUPERIOR — PROBLEMA (Rojos: oscuro → medio → claro)
% ==================================================================

% Nivel 1: Contexto Nacional (más ancho, rojo oscuro)
\fill[red1, rounded corners=4pt,
      blur shadow={shadow blur steps=5, shadow xshift=0.8pt, shadow yshift=-0.8pt}]
    (-9.0, -0.7) rectangle (9.0, -2.5);
\node[text=white, font=\bfseries\large] at (0, -1.1)
    {Contexto Nacional};
\node[text=redTextL, font=\normalsize, text width=16cm] at (0, -1.9)
    {El café aporta el 25,8\% del valor agregado agrícola
     de Colombia \ \ $\bullet$ \ \ 553 M USD en exportaciones
     \ \ $\bullet$ \ \ 540000+ familias cafeteras};

% Nivel 2: El Problema (medio, rojo medio)
\fill[red2, rounded corners=4pt,
      blur shadow={shadow blur steps=5, shadow xshift=0.8pt, shadow yshift=-0.8pt}]
    (-7.3, -3.1) rectangle (7.3, -5.9);
\node[text=white, font=\bfseries\large] at (0, -3.5)
    {El Problema};
\node[text=white, font=\normalsize, text width=12.5cm] at (0, -4.8)
    {El secado solar tradicional ($>$10 días) depende del clima,
     no controla temperatura ni humedad, y genera reacciones
     degenerativas que afectan valor nutricional, textura, color,
     aroma y sabor. Secadores mecánicos convencionales:
     eficiencia $\leq$ 35\%.};

% Nivel 3: Brecha Identificada (más estrecho, rojo claro)
\fill[red3, rounded corners=4pt,
      blur shadow={shadow blur steps=5, shadow xshift=0.8pt, shadow yshift=-0.8pt}]
    (-5.8, -6.5) rectangle (5.8, -8.7);
\node[text=white, font=\bfseries\large] at (0, -6.9)
    {Brecha Identificada};
\node[text=white, font=\normalsize, text width=9.8cm] at (0, -8.0)
    {No existen estudios que evalúen el control
     independiente de la humedad relativa (HR)
     en el secado de café de especialidad
     con bomba de calor modificada.};

% ==================================================================
%  CENTRO — PREGUNTA DE INVESTIGACIÓN (Morado)
% ==================================================================

% Recuadro central morado
\node[
    draw=purple1, line width=2.5pt,
    fill=purplePale,
    rounded corners=8pt,
    inner sep=12pt,
    blur shadow={shadow blur steps=8, shadow xshift=1pt, shadow yshift=-1pt},
    font=\bfseries\large,
    text width=12cm
] (centro) at (0, -10.9) {
    \centering
    \textcolor{purple1}{Pregunta de Investigación}\\[6pt]
    \textcolor{textDark}{\normalfont\normalsize
    ¿Cómo escalar una tecnología de secado de café basada en el control de la humedad relativa mediante bomba de calor para mejorar la calidad y eficiencia del proceso en cafés diferenciados?}
};

% ==================================================================
%  PARTE INFERIOR — IMPACTO (Azules: claro → medio → oscuro)
% ==================================================================

% Nivel 3': Justificación Tecnológica (estrecho, azul claro)
\fill[blue3, rounded corners=4pt,
      blur shadow={shadow blur steps=5, shadow xshift=0.8pt, shadow yshift=-0.8pt}]
    (-5.8, -13.1) rectangle (5.8, -15.3);
\node[text=white, font=\bfseries\large] at (0, -13.5)
    {Justificación Tecnológica};
\node[text=white, font=\normalsize, text width=9.8cm] at (0, -14.5)
    {El SBC modificado consume 60--80\% menos
     energía, permite control preciso de HR, opera independientemente
     del clima y no emite contaminantes.};

% Nivel 2': Justificación Sostenible y Normativa (medio, azul medio)
\fill[blue2, rounded corners=4pt,
      blur shadow={shadow blur steps=5, shadow xshift=0.8pt, shadow yshift=-0.8pt}]
    (-7.3, -15.9) rectangle (7.3, -18.3);

\node[text=white, font=\bfseries\large] at (0, -16.35)
    {Justificación Sostenible y Normativa};

\node[text=blueTextL, font=\normalsize, text width=13.8cm, align=center] at (0, -17.4)
{
Alineación con el Pacto Verde Europeo (2024):
trazabilidad, libre de deforestación y producción sostenible.
Ventaja competitiva para café orgánico colombiano en mercados premium.
Alineación con PIIOM de Minciencias.
};

% Nivel 1': Impacto Esperado (más ancho, azul oscuro)
\fill[blue1, rounded corners=4pt,
      blur shadow={shadow blur steps=5, shadow xshift=0.8pt, shadow yshift=-0.8pt}]
    (-9.0, -18.9) rectangle (9.0, -21.0);
\node[text=white, font=\bfseries\large] at (0, -19.3)
    {Impacto Esperado};
\node[text=blueTextL, font=\normalsize, text width=16cm] at (0, -20.3)
    {Estandarización del secado para café de especialidad
     \ \ $\bullet$ \ \ Mayor competitividad de 22.792+ caficultores del Valle del Cauca
     \ \ $\bullet$ \ \ Posicionar a Colombia como líder en innovación poscosecha sostenible
     \ \ $\bullet$ \ \ Fortalecimiento del capital humano en I+D+i};

% ==================================================================
%  ANOTACIONES LATERALES
% ==================================================================

% Izquierda arriba: PROBLEMA (con fondo blanco para evitar tachado)
\draw[red1, line width=2pt, -{Stealth[length=5pt]}]
    (-9.8, -0.7) -- (-9.8, -8.0);
\node[rotate=90, font=\bfseries\normalsize, text=red1,
      fill=white, inner sep=4pt, rounded corners=2pt] at (-9.8, -4.3)
    {PROBLEMA};

% Izquierda abajo: IMPACTO
\draw[blue1, line width=2pt, {Stealth[length=5pt]}-]
    (-9.8, -13.8) -- (-9.8, -21.0);
\node[rotate=90, font=\bfseries\normalsize, text=blue1,
      fill=white, inner sep=4pt, rounded corners=2pt] at (-9.8, -17.4)
    {IMPACTO};

% Derecha: FOCO
\draw[purple1, line width=1.5pt, dashed]
    (9.8, -9.3) -- (9.8, -12.3);
\node[rotate=-90, font=\bfseries\normalsize, text=purple1,
      fill=white, inner sep=4pt, rounded corners=2pt] at (9.8, -10.8)
    {FOCO};

% ==================================================================
%  SILUETA DECORATIVA DEL RELOJ DE ARENA
% ==================================================================
\draw[grayLine, line width=1.5pt, dashed, opacity=0.45]
    (-9.2, -0.5) -- (-4.8, -10.0) -- (-9.2, -19.2);
\draw[grayLine, line width=1.5pt, dashed, opacity=0.45]
    ( 9.2, -0.5) -- ( 4.8, -10.0) -- ( 9.2, -19.2);

\end{tikzpicture}%
}% fin de \resizebox
\caption{Diagrama de reloj de arena: justificación de la investigación sobre
         secado asistido por bomba de calor para café de especialidad.}
\label{fig:hourglass-justificacion}
\end{figure}
